\documentclass[10pt,a4paper]{amsart}
\usepackage[margin=19mm]{geometry}
\usepackage{amsmath,amssymb,mathtools}
\usepackage[scr=rsfs,cal=euler]{mathalfa}
\usepackage{xfrac}
\usepackage[unicode,hidelinks,bookmarks=false]{hyperref}
\newcommand{\ie}{i.e.\ }
\newcommand{\cf}{cf.\ }
\newcommand{\resp}{respectively\ }
\newcommand{\Subsection}{\subsection}
\providecommand{\nobreakdash}{\nobreak\hbox{-}}
\setlength{\emergencystretch}{2em}
\multlinegap0pt
\usepackage{amssymb}
\usepackage{mathtools}
\usepackage{natbib}
\usepackage{xfrac}
\usepackage{xcolor}
\newtheorem{defi}{Definition}
\newtheorem{prop}{Proposition}
\renewcommand{\H}{\mathcal{H}}
\newcommand{\N}{\mathcal{N}}
\newcommand\bb[1]{\mathbf{#1}}
\newcommand{\bs}[1]{\boldsymbol{#1}}
\newcommand{\bu}{\bb{u}}
\DeclareMathOperator{\curl}{curl}
\newcommand{\ddt}{\frac{d}{dt}}
\renewcommand{\div}{\operatorname{div}}
\renewcommand{\hom}{H^1(\tn)}
\newcommand{\htwo}{H^2(\Omega)}
\newcommand{\io}{\int_{\Omega}}
\newcommand{\ld}{L^2(\Omega)}
\newcommand{\loc}{\operatorname{loc}}
\newcommand{\ot}{(0,\infty)\times\Omega}
\newcommand{\tn}{\Omega}
\newcommand{\ul}{\frac{1}{2}}
\newcommand{\wspnd}{2\mu+\lambda}
\title{\bf Existence, uniqueness, and time-asymptotics of regular solutions in multidimensional thermoelasticity on domains with boundary}
\author{Piotr Micha{\l} Bies}
\date{Source: arXiv:2512.17767v2 (CC BY 4.0)}
\begin{document}
\maketitle

\noindent\emph{Source and licence.} \bf Existence, uniqueness, and time-asymptotics of regular solutions in multidimensional thermoelasticity on domains with boundary. Version arXiv:2512.17767v2 (CC BY 4.0), distributed by arXiv under the Creative Commons Attribution 4.0 International licence, http://creativecommons.org/licenses/by/4.0/, as recorded in the arXiv licence metadata for this exact version and shown on the abstract page https://arxiv.org/abs/2512.17767v2. Full licence notice: \texttt{licenses/arxiv-2512.17767-CC-BY-4.0.txt}.\par\smallskip

\noindent\textit{Editorial scope.} Counted unit: the article's prop twcontdep (source0.tex lines 959--968). The article's complete original proof of it is delivered with it (source0.tex lines 969--1028, one \texttt{proof} environment of 60 lines). No mathematics is written, repaired or re-derived: the statement and the proof are the article's own lines, in the article's own notation, and no retained line is substituted or re-flowed.  Retained uncounted as the source-local context the proof needs: lines 121--128; lines 209--242.  Nothing was omitted from any retained span, so the package carries no omission record.  No retained line contains a citation, so the wrapper carries no bibliography and no bibliography entry.  No retained line contains a figure environment, an image-inclusion command or drawing code, so no image is delivered and the package declares no asset dependency.  Editorial formatting only, in addition: an AMS \texttt{amsart} preamble that restates the article's own declarations and macros used by the retained lines, namely the environments \texttt{defi}, \texttt{prop} on separate counters, together with the standard packages those lines need.  The retained environments print in this document as Definition 1, Proposition 1 in order of appearance, whereas the article prints the same items with the numbering of its own full document.  The complete native source of the article as distributed in the arXiv e-print for this exact version is bundled unchanged as \texttt{sources/191344/source0.tex}.
\begin{equation}\label{system}
\begin{cases}
\bb{u}_{tt}- (2\mu+\lambda)\nabla\div\bb u+\mu\curl\curl\bb u =-\nu\nabla\theta & \quad \text{ in } \ot,\\
\theta_t - \Delta \theta =- \nu\theta \div\bb{u}_t& \quad \text{ in }  \ot,\\
\bb u\cdot\bb n=0,\ \N\bb u=0,\ \nabla\theta\cdot\bb n=0 &\quad\text{ on } (0,\infty)\times\partial\Omega,\\
\bb{u}(0,\cdot)=\bb{u}_0,\ \bb{u}_t(0,\cdot) = \bb{v}_0,\ \theta(0,\cdot)=\theta_0>0,&
\end{cases}
\end{equation}

\begin{defi}\label{soldef}
We say that $(\bb u,\theta)$ is a solution to (\ref{system}) if:
\begin{itemize}
\item The initial data is of regularity
\begin{align*}
\bb u_0\in \htwo\cap\H,\quad \bb v_0\in\H,\quad \theta_0\in\hom\cap L^{\infty}(\tn).
\end{align*}
We also require that there exists $\tilde{\theta}>0$ such that $\theta_0(x)\geq\tilde{\theta}$ for all $x\in\tn$.
\item Solutions $\theta$ and $\bu$ satisfy
\begin{align}\label{reggl}
\begin{split}
\bu&\in L^{\infty}(0,\infty;\htwo\cap\H)\cap W^{1,\infty}(0,\infty;\H)\cap W^{2,\infty}(0,\infty;\ld),\\
\theta&\in L^{\infty}(0,\infty;\hom)\cap L^2_{\loc}(0,\infty;\htwo)\cap H^1_{\loc}(0,\infty; \ld).
\end{split}
\end{align}
\item The momentum equation
\begin{align}\label{defmom}
 \io\bb u_{tt}\cdot\bs\phi dx+(\wspnd)\io\div\bb u\div\bs\phi dx+\mu\io\curl\bb u\cdot\curl\bs\phi dx =-\nu\io\nabla\theta\cdot\bs\phi dx
\end{align}
is satisfied for almost all $t\in (0,\infty)$ and $\bs\phi\in\H$.
\item The entropy equation
\begin{align}\label{defheat}
\io\theta_t\psi dx+\io\nabla\theta\cdot\nabla\psi=-\nu\io\theta \div\bb u_{t}\psi dx
\end{align}
holds for almost all $t\in (0,\infty)$ and for all $\psi\in\hom$.
\item Initial conditions are attained in the following sense:
\[
\theta \in C([0,\infty);\hom),
\]
\[
u\in C([0,\infty);\H),\ u_t \in C([0,\infty);\ld).
\]
\end{itemize}
\end{defi}

\begin{prop}\label{twcontdep}
Let us take two triples of the initial values $\bu_{0,1},\bu_{0,2}\in \H\cap\htwo, \bb v_{0,1},\bb v_{0,2}\in\H$ and $\theta_{0,1},\theta_{0,2}\in\hom\cap L^{\infty}(\tn)$. Moreover, we  require that there exists $\tilde\theta_1>0$ and $\tilde\theta_2>0$ such that $\tilde\theta_1\leq\theta_{0,1}$ and $\tilde\theta_2\leq\theta_2$.
Let us also assume that $(\bu_1,\theta_1)$ is a solution of \eqref{system} starting from $\bu_{0,1},\bb v_{0,1},\theta_{0,1}$ and $(\bu_2,\theta_2)$ is a
solution starting from $\bu_{0,2},\bb v_{0,2},\theta_{0,2}$. Both of the solutions are considered in the sense of Definition \ref{soldef}. Then, for all $t\in(0,\infty)$ there exists the constant $C>0$  such that
\begin{align*}
&\|\bu_1(t,\cdot)-\bu_2(t,\cdot)\|_{\H}+\|\bu_{1,t}(t,\cdot)-\bu_{2,t}(t,\cdot)\|_{\ld}+\|\theta_1(t,\cdot)-\theta_2(t,\cdot)\|_{\ld}\\
&\quad\leq C\left(\|\bu_{0,1}-\bu_{0,2}\|_{\H}+\|\bb v_{0,1}-\bb v_{0,2}\|_{\ld}+\|\theta_{0,1}-\theta_{0,2}\|_{\ld}\right)
\end{align*}
and $C=C(t,d,\mu,\lambda,|\nu|, \|\nabla \bu_{1,t}\|_{L^{\infty}(0,\infty;\ld)},\|\theta_2\|_{L^{\infty}((0,\infty)\times\tn)},\|\nabla \theta_2\|_{L^{\infty}(0,\infty;\ld)}, \|\nabla^2\theta_2\|_{L^2(0,t;\ld)})$.
\end{prop}

\begin{proof}
Let us denote $\bu:=\bu_1-\bu_2$ and $\theta:=\theta_1-\theta_2$. We subtract the equations for $(\bu_2,\theta_2)$ from the equations for $(\bu_1,\theta_1)$ and
obtain
\begin{align}\label{uklpom}
\begin{cases}
\bu_{tt}-(\wspnd)\nabla\div\bu+\mu\curl\curl\bu=-\nu\nabla\theta,\\
\theta_t-\Delta\theta=-\nu\theta \div \bu_{1,t}-\mu \theta_2\div \bu_{t}.
\end{cases}
\end{align}
We multiply by $\bu_t$ the first of the above equations and arrive at
\begin{align}\label{glrow1}
\ul\ddt\left(\io \bu_t^2dx+(\wspnd)\io (\div \bu)^2dx+\mu\io |\curl \bu|^2dx \right)&\leq-\nu\io \bu_t\cdot\nabla\theta dx\nonumber\\
&\leq C\|\bu_t\|_{\ld}^2+\epsilon\|\nabla \theta\|^2_{\ld}.
\end{align}
Whereas, multiplying the second equation in \eqref{uklpom} by $\theta$, we obtain
\begin{align}\label{glrow2}
\nonumber\ul\ddt&\io\theta^2 dx+\io|\nabla\theta|^2 dx=-\nu\io\theta^2\div \bu_{1,t}dx-\nu\io\theta\theta_2\div \bu_{t}dx\\
&=-\nu\io\theta^2\div \bu_{1,t}dx+\nu\io\theta\nabla\theta_2\cdot \bu_{t}dx+\mu\io\theta_2\nabla\theta\cdot \bu_{t}dx=I_1+I_2+I_3.
\end{align}

We bound the integral $I_1$ using the Gagliardo-Nirenberg inequality and then the Young inequality. We get
\begin{align}\label{inconde1}
I_1\leq |\nu|\|\div \bu_{1,t}\|_{L^{\infty}(0,T;\ld)}\|\theta\|_{L^4(\tn)}^2\leq \epsilon\|\nabla\theta\|^2_{\ld}+C\|\theta\|_{\ld}^2.
\end{align}
The integral $I_{3}$ is estimated using the Schwarz inequality and the Young inequality. We also utilize that $\theta_2\in L^{\infty}(\ot)$ for the class of solutions given in Definition \ref{soldef}.
\begin{align}\label{inconde2}
I_{3}\leq |\nu|\|\theta_2\|_{L^{\infty}((0,T)\times\tn)}\|\nabla\theta\|_{\ld}\|\bu_t\|_{\ld}\leq \epsilon\|\nabla\theta\|_{\ld}^2+C\|\bu_t\|_{\ld}^2.
\end{align}

Let us proceed with $I_{2}$. We use the Gagliardo-Nirenberg inequality and the Young inequality.
\begin{align*}
I_{2}&\leq |\nu| \|\theta\|_{L^4(\tn)}\|\nabla\theta_2\|_{L^4(\tn)}\|\bu_t\|_{\ld}\\
&\leq C(\|\nabla\theta\|_{\ld}+\|\theta\|_{\ld})(\|\nabla^2\theta_2\|_{\ld}+\|\nabla\theta_2\|_{\ld})\|\bu_t\|_{\ld}.
\end{align*}
Next, after the application of the Young inequality, we get
\begin{align}\label{inconde3}
I_{2}\leq \epsilon\|\nabla\theta\|_{\ld}^2+ C(\|\nabla^2\theta_2\|_{\ld}^2\|\bu_t\|_{\ld}^2+\|\bu_t\|_{\ld}^2+\|\theta\|_{\ld}^2).
\end{align}

We plug \eqref{inconde1}, \eqref{inconde2} and \eqref{inconde3} into \eqref{glrow2}. The obtained inequality we add to \eqref{glrow1}. It leads us to
\begin{align*}
\ul\ddt&\left(\|\bu_t\|^2_{\ld}+(\wspnd)\|\div \bu\|_{\ld}^2+\mu\|\curl \bu\|_{\ld}^2+\|\theta\|^2_{\ld}\right)+\|\nabla\theta\|_{\ld}^2
\\ &\leq4\epsilon\|\nabla\theta\|_{\ld}^2+ C(\|\nabla^2\theta_2\|_{\ld}^2\|\bu_t\|_{\ld}^2+\|\bu_t\|_{\ld}^2+\|\theta\|_{\ld}^2).
\end{align*}
We take $\epsilon:=\frac18$.

Let us denote $y:=\|\bu_t\|^2_{\ld}+(\wspnd)\|\div \bu\|_{\ld}^2+\mu\|\curl \bu\|_{\ld}^2+\|\theta\|^2_{\ld}$. Then, we rewrite the above inequality as follows
\begin{align*}
y'\leq C y(1+\|\nabla^2\theta_2\|_{\ld}^2).
\end{align*}
We multiply the above inequality by $\operatorname{exp}\left(-C\int_0^s\|\nabla^2\theta_2\|_{\ld}^2dz-Cs\right)$ for $s\in(0,t]$. We obtain
\begin{align*}
\frac{d}{ds}\left(y \operatorname{exp}\left(-C\int_0^s\|\nabla^2\theta_2\|_{\ld}^2dz-Cs\right)\right)\leq 0.
\end{align*}
We integrate over $[0,t]$ and see that
$$
y(t)\leq y(0)\operatorname{exp}\left(C\int_0^t\|\nabla^2\theta_2\|_{\ld}^2ds+Ct\right).
$$
It gives the thesis.
\end{proof}

\end{document}
